% !TeX program = xelatex
\documentclass[11pt]{article}
\usepackage[final]{acl}
\usepackage{fontspec}
\IfFontExistsTF{TeX Gyre Termes}{\setmainfont{TeX Gyre Termes}}{\setmainfont{Times New Roman}}
\usepackage{amsmath,booktabs,tabularx}
\renewcommand{\abstractname}{Tóm tắt}
\renewcommand{\refname}{Tài liệu tham khảo}
\title{Sinh chú thích ảnh tiếng Việt trên KTVIC}
\author{Nguyễn Đức Minh (24022405) \quad Đoàn Quang Mạnh (24022339)\\Nhóm 10 -- Đề tài 05 -- Môn Học sâu}
\begin{document}
\maketitle
\begin{abstract}
Nhóm nghiên cứu bài toán nhận một ảnh RGB và sinh một câu mô tả tiếng Việt. Trên bộ dữ liệu KTVIC, nhóm so sánh CNN encoder kết hợp LSTM có attention với Transformer decoder dùng chung đặc trưng thị giác. Quy trình thống nhất cách chia dữ liệu, từ vựng và đánh giá bằng BLEU, CIDEr; bổ sung phân tích lỗi và chi phí tính toán. Bản đề xuất này xác định phạm vi, thiết kế thí nghiệm và phân công cho hai thành viên, sử dụng lại nền tảng mã nguồn hiện có của dự án.
\end{abstract}

\section{Bài toán và mục tiêu}
Đầu vào là ảnh RGB, đầu ra là một câu mô tả nội dung chính bằng tiếng Việt. Mô hình học phân phối $p(y\mid I)=\prod_t p(y_t\mid y_{<t},I)$. Câu hỏi nghiên cứu là hai loại decoder khác nhau thế nào về chất lượng mô tả, lỗi đối tượng và quan hệ, cũng như thời gian tính toán khi giữ chung encoder và dữ liệu. Không giả định trước Transformer sẽ tốt hơn trên dữ liệu nhỏ.

\section{Dữ liệu và tiền xử lý}
Nhóm chọn KTVIC, bộ dữ liệu mô tả ảnh tiếng Việt thuộc miền đời sống \citep{pham2024ktvic}. Bản dữ liệu local có 3.769 ảnh official train, mỗi ảnh năm chú thích. Sử dụng split hiện có với seed 42: 3.202 ảnh train, 565 validation và 558 official test; hai bản ảnh phía train trùng byte với test đã được cách ly. Các chú thích của một ảnh luôn thuộc cùng một tập.

Văn bản dùng \texttt{segment\_caption}, chuẩn Unicode NFC, chuyển chữ thường, tách theo khoảng trắng và giữ dấu gạch dưới. Vocabulary chỉ học từ train, ngưỡng tần suất 2; độ dài tối đa 64 token gồm BOS/EOS. Ảnh chuyển RGB, chuẩn hóa theo encoder, kích thước $224\times224$. Không dùng test để chọn siêu tham số. Audit cũ ghi nhận 13 ảnh validation gần trùng train; đây là hạn chế phải công bố.

\section{Phương pháp và thí nghiệm}
\textbf{Mô hình bắt buộc A1:} ResNet50 tiền huấn luyện trích xuất lưới $7\times7$; LSTM dùng additive soft attention để tổng hợp vùng ảnh theo từng bước sinh, theo hướng \citet{xu2015attention}.

\textbf{Mô hình so sánh T1:} dùng cùng CNN, chiếu đặc trưng về 256 chiều; Transformer decoder có hai lớp, bốn heads, masked self-attention và cross-attention vào ảnh \citep{vaswani2017attention}. B0 (CNN toàn cục + LSTM không attention) là baseline bổ sung, không thay thế A1.

Đóng băng encoder trong so sánh chính. Huấn luyện decoder bằng teacher forcing và cross-entropy bỏ PAD, giữ EOS; khởi đầu batch 32, tối đa 25 epochs, dừng sớm theo CIDEr validation. Ghi đầy đủ seed, optimizer, số tham số học được, cấu hình và thời gian của từng mô hình.

Đánh giá corpus BLEU-1/2/3/4 và CIDEr với đủ năm references mỗi ảnh, cùng quy tắc xử lý prediction/reference. Chọn checkpoint và greedy hoặc beam độ rộng 3 trên validation; khóa protocol trước test. Ablation dự kiến thay attention động bằng ngữ cảnh trung bình; phân tích thêm lỗi đối tượng, hành động, số lượng, quan hệ, lặp từ và kiểm tra tráo ảnh. Không diễn giải attention map như bằng chứng nhân quả. CLIP là mở rộng nếu đủ thời gian.

\section{Kế hoạch sáu tuần}
\begin{center}
\begin{tabularx}{\columnwidth}{@{}lX@{}}
\toprule
Tuần & Công việc và mốc bàn giao\\
\midrule
1 & Đọc tài liệu, dữ liệu, môi trường; proposal.\\
2 & Baseline, kiểm tra pipeline, log huấn luyện.\\
3 & A1 và T1; milestone tối đa 4 trang.\\
4 & So sánh, ablation, tinh chỉnh trên validation.\\
5 & Mở rộng, phân tích lỗi; nháp báo cáo.\\
6 & Báo cáo ACL, mã nguồn, slide và demo.\\
\bottomrule
\end{tabularx}
\end{center}

\textbf{Phân công dự kiến:} Nguyễn Đức Minh phụ trách dữ liệu, tokenizer, CNN và LSTM attention. Đoàn Quang Mạnh phụ trách Transformer, bộ đánh giá và so sánh thí nghiệm. Cả hai cùng rà mã của nhau, phân tích lỗi, viết báo cáo và chuẩn bị vấn đáp. Phân công này không xác nhận tỷ lệ đóng góp thực tế.

\section{Sản phẩm và trạng thái}
Sản phẩm gồm repo riêng tư với README, mã đọc dữ liệu, kiểm tra tái lập, log, báo cáo ACL và demo. Dữ liệu, trọng số lớn được cung cấp qua nguồn tải riêng. Dùng CPU để kiểm tra pipeline; GPU huấn luyện cần được bố trí và ghi cấu hình thực tế.

Bản này soạn ngày 09/11/2026 từ dự án đã có mã và các kết quả ở giai đoạn sau; lịch sáu tuần là khung học phần, không phải bằng chứng nộp đúng hạn hay khẳng định dự án mới bắt đầu. Repo tuần 1 là bản chọn lọc cho mốc đầu, không công bố lại kết quả test. Đề yêu cầu ba người; nhóm hai người cần được giảng viên chấp thuận. AI hỗ trợ chuẩn bị mã, kiểm tra và soạn thảo; thành viên chịu trách nhiệm kiểm chứng và giải thích. Tài liệu tham khảo dưới đây cần được hai thành viên đọc và đối chiếu.
\bibliography{references}
\end{document}
